Jupyter notebooks are widely used to share computational research, but a
large share of published notebooks fail to re-execute in a different
environment. The FAIR Jupyter project has studied this problem at scale and
built a pipeline that containerises each repository before execution to
remove many such failures; even so, this pipeline only detects and logs the
dependency-related failures that remain -- it does not explain or repair
them. This thesis adds a repair layer that attaches at exactly that point.

The working dataset is a residual set of 214 dependency-related failures
recorded by the Docker-based pipeline, of which 200 fall inside a pip-only
repair scope defined for this first version of the system. Given a failed
notebook, a deterministic classifier decides its subtype and repair
eligibility before any language model is consulted.

An explanation component produces a plain-language, schema-validated
description of the failure for every classified record, independent of
whether it can later be repaired; it has now been run across all 200
repair-scope rows (the 13-row development split and the 187-row evaluation
split), where it produced a schema-valid explanation for 185 of the 187
evaluation records.

A repair component proposes a concrete
fix using an open,
locally-run large language model whose suggestion is grounded in live
metadata from the Python Package Index and validated, field by field,
against that evidence, so that a proposed package name or version is never
invented. A fix-application component applies a validated proposal inside
an environment rebuilt to match the one the notebook originally failed in,
re-executes the whole notebook, and classifies the result as fixed, still
failing, or not meaningfully testable.

All of these components are implemented and covered by 332 automated tests.

Beyond this component-level testing, the complete repair layer was run
once, unattended, over the reserved 187-row evaluation split, using a real
local language model, the live PyPI registry, and a real Docker daemon
throughout. No notebook in this split reached completely successful
re-execution within the two-round repair budget, but this is not the same
as saying no repair worked: 84.9\% of the 73 real repair attempts made
(across both rounds) removed the specific dependency error they targeted,
and every one of the 73 language-model repair responses obtained during
the run produced a schema-valid, PyPI-grounded proposal. The gap between
these two numbers has two main causes. First, the repair agent's static
import-to-distribution mapping only covered enough failing modules to reach
the language model for 73 of 235 repair-agent invocations across both
rounds (31.1\%); the remaining invocations abstained beforehand rather than
guess an unsupported package name. Second, resolving one dependency error
frequently exposed a second, previously hidden one, and the bounded,
two-round repair budget was not always enough to clear all of them within
the same notebook.

A result-logging component joins the classification, explanation, and
repair records these components already produce into one persistent record
per repair attempt, stored in a SQLite benchmark table, and an RDF mapping
re-expresses that record and links it to the existing FAIR Jupyter knowledge
graph. Both were checked directly against real, pilot-scale data: the
mapping produces syntactically valid RDF that links correctly to existing
notebook resources. The benchmark table has since been populated at the
scale of the full 187-row evaluation run; enriching the knowledge graph from
that larger run, rather than from the pilot sample, remains future work.

All 214 dependency-error
records were confirmed to share notebook identifiers with the inspected
knowledge-graph source data.

This thesis reports work completed through a full evaluation of the
repair layer over the reserved 187-row evaluation split, including a
Round-1-versus-final comparison drawn from that same run. Full-scale
enrichment of the knowledge graph from this evaluation run, an
explanation-quality user study, and a baseline comparison against a
non-grounded repair agent remain future work and are stated as such, rather
than presented as already measured.
